\documentclass[11pt]{article}
\usepackage[margin=1in]{geometry}
\usepackage{amsmath,amssymb,amsfonts,mathtools,bm}
\usepackage[T1]{fontenc}
\usepackage[utf8]{inputenc}
\title{Inline and Display Mathematics}
\author{Source-backed grouped formula sample}
\date{}
\begin{document}
\maketitle
\section{Expressions}
\paragraph{Expression 1.} The following expression is taken from a source-backed formula corpus.

\begin{align*}\left| \sum_{\alpha \beta \gamma \delta}\langle \Phi_{\alpha} \Phi_{\beta} \Phi_{\gamma} \Phi_{\delta} \rangle_c\left. \right/ \sum_{\alpha \beta \gamma \delta}\langle \Phi_{\alpha} \Phi_{\beta} \Phi_{\gamma} \Phi_{\delta} \rangle\right|=\left| \langle \Phi^4(x) \rangle_c \left. \right/\langle \Phi^4(x) \rangle \right|\; (n=4) \; ,\end{align*}

\paragraph{Expression 2.} The following expression is taken from a source-backed formula corpus.

\begin{align*}\int^{2\pi}_{0} d\sigma X^{i}(\sigma)\hat{P}_{i}(\sigma) =x^{i}_{0}\hat{p}^{i}_{0} +\frac{g_{ij}}{2\sqrt{\alpha'}} \sum_{n=1}^{\infty}\left\{ \left( \chi_{n}^{i}\alpha^{j}_{-n} +\bar{\chi}^{i}_{n}\tilde{\alpha}^{j}_{-n} \right) +\left(\bar{\chi}_{n}^{i}\alpha^{j}_{n} +\chi_{n}^{i} \tilde{\alpha}^{j}_{n}\right) \right\}~.\end{align*}

\paragraph{Expression 3.} The following expression is taken from a source-backed formula corpus.

\begin{align*}\Biggl( \phi_1(x), {\phi_2(x)\choose\phi_3(x)},{\phi_4(x)\choose\phi_5(x)} \Biggr) \rightarrow \Biggl( c\, e\sp{\lambda x} \phi_1(x), U {e\sp{-\lambda\sp\prime x}\phi_2(x)\choose e\sp{-\lambda\sp\prime x}\phi_3(x)},V {e\sp{\lambda\sp{\prime\prime}x} \phi_4(x)\choose e\sp{\lambda\sp{\prime\prime}x} \phi_5(x)} \Biggr)\end{align*}

\paragraph{Expression 4.} The following expression is taken from a source-backed formula corpus.

\begin{align*}\frac{ -i\langle T a^{\dagger t}_{ {\bf{k}} }(-{\bf{q}})a^{\dagger}_{ {\bf{k}}^{'} }({\bf{q}}^{'}) \rangle }{\langle T 1 \rangle} = \sum_{n}\text{ }exp(\omega_{n}t)\text{ }\frac{ -i\langle T a^{\dagger n}_{ {\bf{k}} }(-{\bf{q}})a^{\dagger}_{ {\bf{k}}^{'} }({\bf{q}}^{'}) \rangle }{\langle T 1 \rangle}\end{align*}

\paragraph{Expression 5.} The following expression is taken from a source-backed formula corpus.

\begin{align*}{\bar G}^{(\zeta )}(\beta )={(\beta +i\pi )\sinh\left( {\displaystyle \pi\beta\over \displaystyle 2\pi +\zeta}\right) \over\Gamma \left( 1- {\displaystyle\pi-i\beta \over \displaystyle 2\pi+\zeta}\right)\Gamma \left( 1- {\displaystyle\pi+i\beta \over \displaystyle 2\pi+\zeta}\right)}\ ,\end{align*}
\end{document}
